
\section{\BookRecordIndex{T0162}行列式线丛}\label{sec:A-2}

设 $X,Y$ 是巴拿赫空间，以 $\mathcal F(X,Y)$ 表示所有\BookRecordIndex{T0128}弗雷德霍姆算子
$D:X\to Y$ 的空间。\BookRecordIndex{T0128}弗雷德霍姆算子 $D\in\mathcal F(X,Y)$ 的\emph{行列式}定义为
一维实向量空间
\[
  \det(D):=\Lambda^{\max}(\ker D)\otimes\Lambda^{\max}(\ker D^*).
\]
下面将证明：当 $D$ 变化时，这些向量空间 $\det(D)$ 拼成 $\mathcal F(X,Y)$ 上的
局部平凡线丛。这里必须留意用于识别不同一维向量空间的同构。
\sourcefootnote{原书第二版勘误（2019 年 11 月 27 日，
\url{https://people.math.ethz.ch/~salamon/PREPRINTS/jholn-erratum.pdf}）指出，
印刷页 534--546 关于\BookRecordIndex{T0162}行列式线丛的讨论须重写以修正符号。以下保留锁定原书的论述；
涉及定向和同构的符号结论，不应视作已按勘误修正。}

把实直线 $\mathbb R$ 看作一维实向量空间。对任意两个一维实向量空间 $V,W$，
我们用 $V\cong W$ 表示它们\emph{自然同构}，即两者之间有一个显然的同构选择。
例如，当 $V$ 是一维空间时，有自然同构
\[
  V\otimes V^*\longrightarrow\mathbb R:
  v\otimes v^*\longmapsto v^*(v).
\]

自然同构的概念可以用范畴论语言表述得更精确。以 $\mathcal V$ 表示由一维实向量空间
及其同构组成的范畴，令 $\mathcal C$ 为任意其他范畴。两个函子
$\mathcal F_0,\mathcal F_1:\mathcal C\to\mathcal V$ 称为\emph{自然同构}，
如果存在自然变换
$\mathcal T:\operatorname{Ob}(\mathcal C)\to\operatorname{Mor}(\mathcal V)$，
它给每个对象 $A\in\operatorname{Ob}(\mathcal C)$ 指派一个向量空间同构
$\mathcal T(A):\mathcal F_0(A)\to\mathcal F_1(A)$，并且对
$A,B\in\operatorname{Ob}(\mathcal C)$、$T\in\operatorname{Mor}(A,B)$ 满足
\[
  \mathcal T(B)\circ\mathcal F_0(T)
  =\mathcal F_1(T)\circ\mathcal T(A).
\]
许多情形下，$\mathcal T(A)$ 的选择由上下文即可确定，我们便直接写作
$\mathcal F_0(A)\cong\mathcal F_1(A)$。上面的同构记作
$V\otimes V^*\cong\mathbb R$ 就是一个例子。

有限维实向量空间 $V$ 的\emph{最高外幂} $\Lambda^{\max}V$，是 $V$ 中有序
$n$ 元组 $v_1\wedge\cdots\wedge v_n$ 的等价类所成的空间，其中 $\dim V=n$。
两个这样的 $n$ 元组 $v_1\wedge\cdots\wedge v_n$ 和
$w_1\wedge\cdots\wedge w_n$ 等价，当且仅当它们都构成基且所诱导的
$V$ 的同构行列式为 $1$，或者两个元组都线性相关。因此，
$\Lambda^{\max}V$ 中的非零向量确定 $V$ 的一个定向。当 $\dim V=0$ 时，
定义 $\Lambda^{\max}V:=\mathbb R$。

两个一维实向量空间 $V,W$ 的\emph{张量积} $V\otimes W$，是有序对
$(v,w)\in V\times W$ 的等价类 $v\otimes w$ 所成的空间，其中对每个
$\lambda\in\mathbb R$ 都有 $\lambda v\otimes w=v\otimes\lambda w$。
对任意两个有限维实向量空间 $V,W$，有自然同构
\[
  \Lambda^{\max}V\otimes\Lambda^{\max}W
  \cong\Lambda^{\max}(V\oplus W).
\]
（请证明这一点！）

对每个 $n$ 维实向量空间 $V$，把最高外幂
$\Lambda^{\max}V^*:=\Lambda^{\max}(V^*)$ 看作 $V$ 上的交替 $n$-形式空间。
自然同构
\[
  \Lambda^{\max}V\otimes\Lambda^{\max}V^*\cong\mathbb R
\]
由如下映射给出：
\[
  (v_1\wedge\cdots\wedge v_n)
  \otimes(v_1^*\wedge\cdots\wedge v_n^*)
  \longmapsto\det\bigl(v_k^*(v_j)\bigr).
\]
它诱导同构
$\Lambda^{\max}V^*\cong(\Lambda^{\max}V)^*$。
同样，任意同构 $\Lambda^{\max}V\otimes\Lambda^{\max}W\to\mathbb R$
都诱导同构 $\Lambda^{\max}W\to(\Lambda^{\max}V)^*
\cong\Lambda^{\max}V^*$。

\begin{lemma}\label{lem:A-2-1}
设 $V$ 是有限维实向量空间，$W\subset V$ 是线性子空间。则有自然同构
\[
  \Lambda^{\max}V\cong\Lambda^{\max}W\otimes\Lambda^{\max}(V/W).
\]
\end{lemma}

\begin{proof}
令 $N:=\dim V$、$n:=\dim W$。以 $\mathcal B(V,W)$ 表示 $V$ 的所有基
$v_1,\ldots,v_N$ 构成的集合，其中前 $n$ 个向量张成 $W$。
对 $v\in V$，记 $[v]:=v+W$。映射
$\mathcal B(V,W)\to\Lambda^{\max}W\otimes\Lambda^{\max}(V/W)$，
\[
  (v_1,\ldots,v_N)\longmapsto
  (v_1\wedge\cdots\wedge v_n)
  \otimes([v_{n+1}]\wedge\cdots\wedge[v_N])
\]
诱导所需同构。验证这些同构构成自然变换，留给读者。
\end{proof}

经过这些准备，再考虑\BookRecordIndex{T0128}弗雷德霍姆算子族 $D$ 的\BookRecordIndex{T0162}行列式线
\[
  \det(D):=\Lambda^{\max}(\ker D)\otimes\Lambda^{\max}(\ker D^*).
\]

\begin{theorem}\label{thm:A-2-2}
空间
\[
  \det(X,Y):=\bigl\{(D,\sigma)\mid D\in\mathcal F(X,Y),\;
  \sigma\in\det(D)\bigr\}
\]
是 $\mathcal F(X,Y)$ 上的实线丛。
\end{theorem}

\begin{proof}
必须证明 $\det(X,Y)$ 在每个\BookRecordIndex{T0128}弗雷德霍姆算子 $D\in\mathcal F(X,Y)$ 的邻域中
有局部平凡化。先假定 $D$ 是满射。这时存在估计
$\|y^*\|_{Y^*}\le c\|D^*y^*\|_{X^*}$。
因此，只要 $P$ 足够小，$D+P:X\to Y$ 仍是满射。由定理~\ref{thm:A-1-5}，
$D$ 与 $D+P$ 的\BookRecordIndex{T0129}弗雷德霍姆指标相同，故二者的核同维。
选取 $D$ 的一个\BookRecordIndex{T0392}右逆 $T:Y\to X$，使 $DT=\operatorname{id}_Y$。
只要 $P$ 足够小，便有同构
\[
  \ker(D+P)\longrightarrow\ker D:
  x\longmapsto x+TPx.
\]
因此，在满射\BookRecordIndex{T0128}弗雷德霍姆算子的一个邻域中，$D+P$ 的核拼成
$\mathcal F(X,Y)$ 上的局部平凡向量丛。

现在把一般情形化为满射情形。首先，对任意\BookRecordIndex{T0128}弗雷德霍姆算子 $D_0:X\to Y$，
存在非负整数 $N$\sourcefootnote{原书称 $N$ 为正整数，随即又取
$N=\dim\operatorname{coker}D_0$；当 $D_0$ 满射时该维数为 $0$，
故这里据紧邻的构造改为“非负整数”。} 和单射线性映射
$\Phi:\mathbb R^N\to Y$，使得由
\[
  D_0\oplus\Phi:X\oplus\mathbb R^N\longrightarrow Y,
  \qquad (D_0\oplus\Phi)(x,\zeta):=D_0x+\Phi\zeta
\]
定义的算子 $D_0\oplus\Phi$ 为满射。事实上，令
$N=\dim\operatorname{coker}D_0$，并选取 $y_1,\ldots,y_N\in Y$，
使它们张成 $\operatorname{im}D_0$ 在 $Y$ 中的一个补空间。
则线性映射 $\Phi\zeta=\sum_j\zeta_jy_j$ 满足要求。

现在设 $D:X\to Y$ 是\BookRecordIndex{T0128}弗雷德霍姆算子，$\Phi:\mathbb R^N\to Y$ 是使
$D\oplus\Phi$ 为满射的线性映射。我们证明有自然同构
$\det(D)\cong\det(D\oplus\Phi)$。为此，考虑正合序列
\[
  0\longrightarrow\ker D\longrightarrow\ker(D\oplus\Phi)
  \longrightarrow\Phi^{-1}(\operatorname{im}D)\longrightarrow0,
\]
其中第二个映射是 $\ker D\to\ker(D\oplus\Phi):x\mapsto(x,0)$，
第三个映射是
$\ker(D\oplus\Phi)\to\Phi^{-1}(\operatorname{im}D):(x,\zeta)\mapsto\zeta$。
由正合性，后一个映射下降为同构
$\ker(D\oplus\Phi)/\ker D\cong\Phi^{-1}(\operatorname{im}D)$。
因此，由引理~\ref{lem:A-2-1} 得
\begin{equation}\label{eq:A-2-1}
  \Lambda^{\max}\ker(D\oplus\Phi)
  \cong\Lambda^{\max}(\ker D)
  \otimes\Lambda^{\max}\bigl(\Phi^{-1}(\operatorname{im}D)\bigr).
\end{equation}
由引理~\ref{lem:A-2-1} 还可知，对任意线性子空间
$W\subset\mathbb R^N$，有
$\Lambda^{\max}W\otimes\Lambda^{\max}(\mathbb R^N/W)
\cong\Lambda^{\max}\mathbb R^N\cong\mathbb R$，从而
$\Lambda^{\max}W\cong\Lambda^{\max}(\mathbb R^N/W)^*$。
取 $W=\Phi^{-1}(\operatorname{im}D)$，得到
\begin{equation}\label{eq:A-2-2}
  \Lambda^{\max}\bigl(\Phi^{-1}(\operatorname{im}D)\bigr)
  \cong\Lambda^{\max}
  \left(\frac{\mathbb R^N}{\Phi^{-1}(\operatorname{im}D)}\right)^*.
\end{equation}
由于 $Y=\operatorname{im}D+\operatorname{im}\Phi$，有同构
\[
  \frac{\mathbb R^N}{\Phi^{-1}(\operatorname{im}D)}
  \longrightarrow\frac{Y}{\operatorname{im}D}:
  [\zeta]\longmapsto[\Phi\zeta].
\]
于是
\begin{equation}\label{eq:A-2-3}
  \Lambda^{\max}
  \left(\frac{\mathbb R^N}{\Phi^{-1}(\operatorname{im}D)}\right)^*
  \cong\Lambda^{\max}\left(\frac{Y}{\operatorname{im}D}\right)^*
  \cong\Lambda^{\max}(\ker D^*).
\end{equation}
这里用到了商空间 $Y/Y_1$ 的对偶空间与 $Y^*$ 中 $Y_1$ 的零化子同构这一事实。
当 $Y_1:=\operatorname{im}D$ 时，该零化子等于 $\ker D^*$。
由式~\eqref{eq:A-2-2} 和~\eqref{eq:A-2-3} 得到
$\Lambda^{\max}(\Phi^{-1}(\operatorname{im}D))
\cong\Lambda^{\max}(\ker D^*)$。
再结合式~\eqref{eq:A-2-1}，得到
\[
  \det(D)=\Lambda^{\max}(\ker D)\otimes\Lambda^{\max}(\ker D^*)
  \cong\Lambda^{\max}\ker(D\oplus\Phi)=\det(D\oplus\Phi).
\]
这证明了定理~\ref{thm:A-2-2}。
\end{proof}

\begin{exercise}\label{ex:A-2-3}
设 $D:X\to Y$ 是\BookRecordIndex{T0128}弗雷德霍姆算子，$\Phi:\mathbb R^N\to Y$ 是线性映射
（不必为单射），且 $D\oplus\Phi$ 是满射。
\begin{enumerate}[label=(\roman*)]
\item 假设 $0<k:=\dim\ker D$，并且
$0<\ell:=\dim\ker D^*<N$。给定 $\ker D$ 的基
$x_1,\ldots,x_k$ 及 $\ker D^*$ 的基
$y_1^*,\ldots,y_\ell^*$，证明：存在 $\mathbb R^N$ 的基
$\zeta_1,\ldots,\zeta_N$ 和向量
$\xi_1,\ldots,\xi_{N-\ell}\in X$，使得
\[
\begin{aligned}
  \langle y_i^*,\Phi\zeta_{N-\ell+j}\rangle
    &=\delta_{ij}, &&i,j=1,\ldots,\ell,\\
  D\xi_j+\Phi\zeta_j&=0, &&j=1,\ldots,N-\ell,\\
  \det(\zeta_1\cdots\zeta_N)&=1.&&
\end{aligned}
\]
证明如下映射是良定义的线性同构，并与定理~\ref{thm:A-2-2} 的证明中构造的同构一致：
\[
\begin{aligned}
  \Lambda^{\max}(\ker D)\otimes\Lambda^{\max}(\ker D^*)
   &\longrightarrow\Lambda^{\max}\ker(D\oplus\Phi),\\
  (x_1\wedge\cdots\wedge x_k)
  \otimes(y_1^*\wedge\cdots\wedge y_\ell^*)
   &\longmapsto(x_1,0)\wedge\cdots\wedge(x_k,0)\\
   &\hphantom{\longmapsto{}}\wedge(\xi_1,\zeta_1)
   \wedge\cdots\wedge(\xi_{N-\ell},\zeta_{N-\ell}).
\end{aligned}
\]
\item 考察定理~\ref{thm:A-2-2} 中的同构
$\det(D)\cong\det(D\oplus\Phi)$，分别讨论 $k=0$、$\ell=0$，
以及 $k$ 任意而 $\ell=N$ 的情形。
\end{enumerate}
\end{exercise}

\noindent\textbf{交叉数。}
为了更好地理解指标为零时的线丛
$\det(X,Y)\to\mathcal F(X,Y)$，我们将沿道路
$D:[0,1]\to\mathcal F(X,Y)$ 的平凡化解释为交叉数。
记
\[
  \mathcal F^0(X,Y):=\bigl\{D\in\mathcal F(X,Y)
  \mid\operatorname{index}D=0\bigr\}
\]
并对每个整数 $k\ge0$ 考虑指标为零且核为 $k$ 维的\BookRecordIndex{T0128}弗雷德霍姆算子子集
\[
  \mathcal F_k^0(X,Y):=\bigl\{D\in\mathcal F(X,Y)
  \mid\operatorname{index}D=0,\ \dim\ker D=k\bigr\}.
\]
我们断言，这是余维数为 $k^2$ 的子流形。具体地，在
$D\in\mathcal F_k^0(X,Y)$ 处的切空间为
\[
  T_D\mathcal F_k^0(X,Y)
  =\bigl\{P\in\mathcal L(X,Y)\mid
  Px\in\operatorname{im}D\text{ 对每个 }x\in\ker D\bigr\}.
\]
因而，它有一个 $k^2$ 维的补空间，由从 $\ker D$ 到
$\operatorname{im}D$ 在 $Y$ 中的一个补空间的线性算子组成。并集
\sourcefootnote{原书将下式左侧记作 $\mathcal F_1^0(X,Y)$，但上文已用同一记号定义
$k=1$ 的单一\BookRecordIndex{T0072}层；此处保留原记号，并以右侧并集明确其所指。
下文“简单交叉”处则按前一定义理解为核维数等于 $1$。}
\[
  \mathcal F_1^0(X,Y)
  =\bigcup_{k\ge1}\mathcal F_k^0(X,Y)
\]
是一种余维数为 $1$ 的分\BookRecordIndex{T0072}层子簇；它在指标为零的\BookRecordIndex{T0128}弗雷德霍姆算子空间中的补集，
就是可逆算子空间。

现在考虑道路 $[0,1]\to\mathcal F^0(X,Y):t\mapsto D_t$，其端点
$D_0,D_1$ 都可逆。假设这条道路在强算子拓扑下连续可微，并对每个 $x\in X$ 定义
算子 $\dot D_t:X\to Y$ 为
\[
  \dot D_tx:=\frac{d}{dt}D_tx.
\]
若 $\dim\ker D_t>0$\sourcefootnote{原书写作 $\ker D_t>0$，
将向量空间与数比较；据本页的分\BookRecordIndex{T0072}层定义改写为核维数大于零。}，
则称 $t\in[0,1]$ 为一个\emph{交叉点}。若
\[
  x\in\ker D_t,\quad \dot D_tx\in\operatorname{im}D_t
  \quad\Longrightarrow\quad x=0,
\]
则称这个交叉是\emph{正则的}。这意味着 $\dot D_t$ 把 $\ker D_t$
双射到 $\operatorname{im}D_t$ 的一个补空间。若正则交叉 $t$ 还满足
$D_t\in\mathcal F_1^0(X,Y)$，则称之为\emph{简单交叉}。
注意，只要 $t$ 是正则交叉，在充分小的非零 $s$ 下 $D_{t+s}$ 都可逆。
因此，每个正则交叉都是孤立的。若道路 $t\mapsto D_t$ 只有正则交叉，
我们定义其\emph{交叉指标}为
\begin{equation}\label{eq:A-2-4}
  \mu(\{D_t\}):=\sum_t\dim\ker D_t.
\end{equation}

我们将证明，此数模 $2$ 的值是同伦不变量，并决定了沿道路
$t\mapsto D_t$ 平凡化\BookRecordIndex{T0162}行列式线丛 $\det(X,Y)$ 所得到的映射
$\det(D_0)\to\det(D_1)$ 的符号。更具体地，考虑单位区间上的线丛
\[
  L:=\bigl\{(t,\sigma)\mid t\in[0,1],\ \sigma\in\det(D_t)\bigr\}.
\]
该线丛的平凡化诱导一个同构 $\det(D_0)\to\det(D_1)$。
由于一维向量空间 $\det(D_t)$ 从 $X,Y$ 继承范数，这个同构可以唯一地选为等距映射。
又因 $\det(D_0)=\det(D_1)=\mathbb R$，这个同构就是乘以某个绝对值为 $1$ 的实数，
记作 $\nu(\{D_t\})\in\{\pm1\}$。

\begin{proposition}\label{prop:A-2-4}
设 $[0,1]\to\mathcal F^0(X,Y):t\mapsto D_t$ 是连续可微道路，
端点 $D_0,D_1$ 可逆，且仅有正则交叉。则沿这条道路的\BookRecordIndex{T0162}行列式线丛的任意平凡化，
都诱导符号为
\begin{equation}\label{eq:A-2-5}
  \nu(\{D_t\})=(-1)^{\mu(\{D_t\})}
  =\prod_t(-1)^{\dim\ker D_t}
\end{equation}
的同构 $\det(D_0)=\mathbb R\to\det(D_1)=\mathbb R$。
乘积取遍所有交叉点 $t$。特别地，交叉指标模 $2$ 在固定端点的同伦下不变。
\end{proposition}

\begin{proof}
考虑一条在 $t=0$ 处只有一个交叉的道路 $t\mapsto D_t$。
适当选取 $X,Y$ 的分解后，不失一般性，可设 $X=Y=X_0\oplus\mathbb R^k$，且
\[
  D_0=\begin{pmatrix}\operatorname{id}&0\\0&0\end{pmatrix},
  \qquad
  \dot D_0=\begin{pmatrix}A&0\\B&\operatorname{id}\end{pmatrix}.
\]
考虑线性映射 $\Phi:\mathbb R^k\to Y=X_0\oplus\mathbb R^k$，
$\Phi z=(0,z)$。则
\sourcefootnote{原书只假设道路连续可微，却在下式及后续的核嵌入式写 $O(t^2)$；
单由该假设一般只能\BookRecordIndex{T0326}推出 $o(t)$。此处保留原式，未替原证明补加更强\BookRecordIndex{T0408}正则性假设。}
\[
  D_t\oplus\Phi
  =\begin{pmatrix}
    \operatorname{id}+tA&0&0\\
    tB&t\operatorname{id}&\operatorname{id}
   \end{pmatrix}+O(t^2).
\]
因此，算子 $D_t\oplus\Phi$ 的核可由如下形式的嵌入
$\iota_t:\mathbb R^k\to X\oplus\mathbb R^k
=(X_0\oplus\mathbb R^k)\oplus\mathbb R^k$ 平凡化：
\[
  \iota_t=\begin{pmatrix}0\\\operatorname{id}\\-t\operatorname{id}\end{pmatrix}
  +O(t^2).
\]
这里，上面两个分块表示 $\ker(D_t\oplus\Phi)\subset X\oplus\mathbb R^k$
中的 $X$ 分量，第三个分块表示 $\mathbb R^k$ 分量。
对充分小的 $\varepsilon>0$，考虑复合映射
\[
  \mathbb R^k\longrightarrow\ker(D_{-\varepsilon}\oplus\Phi)
  \longrightarrow\ker(D_\varepsilon\oplus\Phi)
  \longrightarrow\mathbb R^k,
\]
其中最后一个同构由投影到 $\mathbb R^k$ 诱导，第一个同构是它的逆，
第二个同构由平凡化给出。由于 $\iota_t$ 中有因子
$-t\operatorname{id}$，此复合映射形如
$-\operatorname{id}+O(\varepsilon)$。当 $k$ 为奇数时它反转定向，
当 $k$ 为偶数时它保持定向。这在仅有一个正则交叉的情形下证明了
式~\eqref{eq:A-2-5}；一般情形显然随之得到。
\end{proof}

\begin{exercise}\label{ex:A-2-5}
设 $X=Y=H$ 是无限维希尔伯特空间。构造一条指标为零、交叉数为一的
\BookRecordIndex{T0128}弗雷德霍姆算子环路。由此\BookRecordIndex{T0326}推出，$H$ 上指标为零的\BookRecordIndex{T0128}弗雷德霍姆算子空间之上的
行列式丛没有平凡化。提示：选取一组正交归一基 $e_0,e_1,e_2,\ldots$。
当 $0\le t\le1/2$ 时，使 $A_t\in\mathcal L(H)$ 在每个
$(e_{2j},e_{2j+1})$ 平面上转动角度 $2\pi t$（$j\ge0$），从而
$A_0=\operatorname{id}$ 且 $A_{1/2}=-\operatorname{id}$。
当 $1/2\le t\le1$ 时，令 $A_te_0=-e_0$，并使 $A_t$ 在每个
$(e_{2j-1},e_{2j})$ 平面上转动角度 $2\pi t$（$j\ge1$）。
此时 $A_1$ 在 $\mathbb Re_0$ 上作用为负恒等映射，在其正交补上作用为恒等映射。
最后，用直线段将 $A_1$ 连到恒等映射。
\end{exercise}

\begin{exercise}\label{ex:A-2-6}
构造一条指标为 $1$ 的满射\BookRecordIndex{T0128}弗雷德霍姆算子环路，使其核组成莫比乌斯带。
\end{exercise}

\begin{exercise}\label{ex:A-2-7}
设 $X,Y$ 是复巴拿赫空间，
$[0,1]\to\mathcal F(X,Y):t\mapsto D_t$ 是一条由复线性\BookRecordIndex{T0128}弗雷德霍姆算子
组成的道路。那么，一维实向量空间 $\det(D_t)$ 从 $\ker D_t$ 与
$\ker D_t^*$ 的复结构继承自然定向。证明：实线丛
$\bigcup_t\det(D_t)$ 的任意平凡化都诱导保持定向的同构
\[
  \det(D_0)\longrightarrow\det(D_1).
\]
提示：考虑复线丛
\[
  \det^c(X,Y)\longrightarrow\mathcal F^c(X,Y),
\]
它在复线性\BookRecordIndex{T0128}弗雷德霍姆算子 $D\in\mathcal F^c(X,Y)$ 处的纤维，
是一维复向量空间
\[
  \det^c(D):=\Lambda_{\mathbb C}^{\max}\ker D
  \otimes_{\mathbb C}\Lambda_{\mathbb C}^{\max}\ker D^*.
\]
这里 $D^*:Y^*\to X^*$ 表示实对偶算子；若 $J:X\to X$ 是乘以 $i$，
则 $X^*$ 上的复结构由对偶算子 $J^*:X^*\to X^*$ 给出。
证明定理~\ref{thm:A-2-2} 的论证可移至复范畴，因此
$\det^c(X,Y)$ 是 $\mathcal F^c(X,Y)$ 上的局部平凡复线丛。
现在，对每个 $n$ 维复向量空间 $V$，有自然的二次映射
\[
  \Lambda_{\mathbb C}^{\max}V\longrightarrow\Lambda_{\mathbb R}^{\max}V:
  \tau\longmapsto i^n(-1)^{n(n-1)/2}\tau\wedge\bar\tau,
\]
其中 $\bar\tau\in\Lambda_{\mathbb C}^{\max}\bar V$。
这里 $\bar V$ 表示赋予相反复结构的向量空间，而
$\Lambda_{\mathbb C}^{\max}V$ 与
$\Lambda_{\mathbb C}^{\max}\bar V$ 都是
$\Lambda_{\mathbb R}^{\mathrm{mid}}V\otimes\mathbb C$
的自然线性子空间。按这些约定，上面的映射将
$v_1\wedge\cdots\wedge v_n$ 送到
$v_1\wedge(Jv_1)\wedge\cdots\wedge v_n\wedge(Jv_n)$。
\end{exercise}
